\section{把证明约化到拉格朗日情形}
\label{sec:app-c-reduction-lagrangian}

现在回忆第 20.2.2 节的设置。

设 $(\Sigma,j)$ 是带三个标记点的亏格 $0$ 黎曼曲面。以 $D_i$（$i=1,2,3$）表示
这三点的邻域。假设 $i=1,2$ 所对应的穿孔为正穿孔，其解析坐标为
\[
 \varphi_i^+\colon D_i\setminus\{z_i\}
 \longrightarrow(-\infty,0]\times S^1,
 \qquad i=1,2,
\]
而第三个穿孔为负穿孔，其解析坐标为
\[
 \varphi_3^-\colon D_3\longrightarrow[0,\infty)\times S^1.
\]
\footnote{译者注：原书在后一穿孔处写“$i=0$”，但同一句的邻域、坐标和下文均为
下标 $3$；译文按第三个穿孔表述。原书还没有在 $D_3$ 中删去 $z_3$，此处照录。}
以 $(\tau,t)$ 表示柱面
\[
 [0,\infty)\times S^1=[0,\infty)\times\mathbb R/\mathbb Z
\]
上的标准柱面坐标，并如前固定截断函数
$\rho^+\colon(-\infty,0]\to[0,1]$ 与
$\rho^-\colon[0,\infty)\to[0,1]$，满足
$\rho^-(\tau)=\rho^+(-\tau)$。

对 $i=1,2$，令 $[z_i,w_i]$ 是 $\mathcal A_{H_i}$ 的临界点，并记
$\widetilde z_i=[z_i,w_i]$，以及
$\widetilde z=(\widetilde z_1,\widetilde z_2,\widetilde z_3)$。取
$u\in\mathcal M(H,\widetilde J;\widetilde z)$，其中
$\mathcal M(H,\widetilde J;\widetilde z)$ 是第 20.2.2 节给出的、与该三元组相伴的
模空间。重述命题 20.2.13 如下。

\begin{proposition}
\label{prop:app-c-pants-dimension}
$\mathcal M(H,\widetilde J;\widetilde z)$ 的（虚）维数为
\begin{align*}
 \dim\mathcal M(H,\widetilde J;\widetilde z)
 &=2n-(-\mu_{H_1}(\widetilde z_1)+n)
       -(-\mu_{H_2}(\widetilde z_2)+n)
       -(\mu_{H_3}(\widetilde z_3)+n)\\
 &=-n+\bigl(\mu_{H_1}(\widetilde z_1)
            +\mu_{H_2}(\widetilde z_2)
            -\mu_{H_3}(\widetilde z_3)\bigr).
\end{align*}
\end{proposition}

只需计算 $u\in\mathcal M(H,\widetilde J;\widetilde z)$ 处线性化算子的弗雷德霍姆
指标。随后用各 $w_i$ 为 $u$ 加帽，并应用粘合公式（或切除公式），即可得到该公式。

我们将证明：当 $z$ 是帽中的\emph{入端}渐近极限时，每个帽的指标为 $-\mu_*(z)+n$。注意，
粘合后的线性化算子是在亏格 $0$ 闭黎曼曲面上、作用于\emph{平凡丛}
$E=\Sigma\times\mathbb R^{2n}$ 的截面空间的扰动柯西--黎曼算子。平凡性来自同伦条件
\[
 [u]\amalg_{i=1}^3[w_i]=0.
\]
粘合算子可经由一族弗雷德霍姆算子光滑同伦到作用在该平凡丛截面空间上的标准多尔博
算子。由于平凡丛的第一陈数 $c_1(E)$ 为零，由指标公式 (10.4.14)，后一算子的
弗雷德霍姆指标为 $2n$。

由弗雷德霍姆指标的粘合公式，
\[
 \operatorname{Index}D_u\bar\partial_{(H,\widetilde J)}
 +\sum_{i=1}^3\operatorname{Index}D_{w_i}\bar\partial_{(H_i,J_i)}=2n,
\]
从而
\begin{equation}
\label{eq:app-c-gluing-index}
 \operatorname{Index}D_u\bar\partial_{(H,\widetilde J)}
 =2n-\sum_{i=1}^3\operatorname{Index}D_{w_i}\bar\partial_{(H_i,J_i)}.
\end{equation}
剩下只须计算 $\operatorname{Index}D_{w_i}\bar\partial_{(H_i,J_i)}$。

现在为每个帽配上一个典范柯西--黎曼型算子，它作用在空间
\[
 W^{1,p}(C_+;\mathbb R^{2n})
\]
上，其中曲面 $C_+$ 为
\[
 C_+=(-\infty,0]\times S^1\cup S_+^2/{\sim}.
\]
这里 $S_+^2$ 是右半球；等价关系 $\sim$ 把
$(0,t)\in(-\infty,0]\times S^1$ 与
$(r,\theta)=(1,t)\in S_+^2$ 等同，其中 $(r,\theta)$ 是
$S_+^2\cong D^2$ 的极坐标。关于 $(-\infty,0]\times S^1$ 上的柱面坐标
$(\tau,t)$，算子 $D_{w_i}\bar\partial_{(H_i,J_i)}$ 具有形式
\[
 \bar\partial_{J,T}=\frac\partial{\partial\tau}
 +J\frac\partial{\partial t}+T,
\]
其中对某个 $R>0$ 和所有 $\tau\leq-R$，
$T(\tau,\cdot)\equiv T(-\infty,\cdot)$ 且
$J(\tau,\cdot)\equiv J(-\infty,\cdot)$。

在不改变弗雷德霍姆指标的情况下形变 $(-\infty,0]\times S^1$ 上的算子，还可假设
对充分大的 $R>0$，
\[
 J(\tau,\cdot)=
 \begin{cases}
 J(\infty,\cdot),&\tau\geq R+1,\\
 i,&\tau\leq R,
 \end{cases}
 \qquad
 T(\tau,\cdot)=
 \begin{cases}
 T(\infty,\cdot),&\tau\geq R+1,\\
 0,&\tau\leq R.
 \end{cases}
\]
\footnote{译者注：原书刚把柱面固定为 $(-\infty,0]\times S^1$，此处却以
$\tau\geq R+1$ 和正无穷端规定形变；两套定义域不相容。译文照录，不自行改号。}
于是该算子典范地延拓到 $S_+^2$ 区域，成为作用在
$W^{1,p}(C_+;\mathbb R^{2n})$ 上的算子。

\begin{theorem}
\label{thm:app-c-cap-index}
设 $T,J$ 如上，并令 $\Psi$ 是方程
\begin{equation}
\label{eq:app-c-cap-fundamental-solution}
 \begin{cases}
 \partial\Psi/\partial t-J(\infty,t)T(\infty,t)\Psi=0,\\
 \Psi(0)=\operatorname{id}
 \end{cases}
\end{equation}
在 $C_+$ 上的解。则
\[
 \operatorname{Index}\bar\partial_{J,T}
 =-\mu_{\mathrm{CZ}}(\Psi)+n,
\]
其中 $\Psi$ 是 (C.4) 的基本解，而 $\mu_{\mathrm{CZ}}(\Psi)$ 是辛道路
$t\mapsto\Psi(t)\in\operatorname{Sp}(2n)$ 的\personRobbin{}--\personSalamon{}广义康利--曾德尔指标。
\footnote{译者注：本定理本身把基本解定义在式 (C.2)，末句却前引尚未出现的式
(C.4)；译文保留原交叉引用。}
\end{theorem}
